% !TeX root = Master.tex
\section{Framework}
\label{sec:framework}
\subsection{Causal framework}
\label{sec:causal-framework}
Let \(X_i\), \(A_i\in\{0,1\}\), and \(Y_i\) denote the covariate vector, binary treatment, and observed outcome for unit \(i\), and let \(Y_i(0)\) and \(Y_i(1)\) denote its potential outcomes under control and treatment, respectively. When no unit index is shown, \(X\), \(A\), \(Y(0)\), \(Y(1)\), and \(Y\) denote the corresponding random variables for a generic unit. In an EMCS in causal inference, the researcher starts with an empirically observed reference dataset assumed to have been generated by a full causal law \(P_0\) for these unit-level variables:
\begin{equation}
\begin{gathered}
\mathcal D_0=\{(X_i,A_i,Y_i)\}_{i=1}^{n_0},\\
(X_i,A_i,Y_i(0),Y_i(1),Y_i)\sim P_0,\qquad i=1,\dots,n_0,
\end{gathered}
\label{eq:observed-data-causal-law}
\end{equation}
where consistency links the observed outcome to the potential outcomes:
\begin{equation}
\text{\textbf{Consistency:}}\quad Y=AY(1)+(1-A)Y(0).
\label{eq:consistency}
\end{equation}

We suppose that the researcher is interested in an estimand that is well defined and can be expressed by averaging the conditional expectation of the \emph{causal contrast}, \(Y(1)-Y(0)\), over a target covariate distribution \(P_X^{\mathrm{tar}}\):
\begin{equation}
\Phi_0(P_X^{\mathrm{tar}})
=
\mathbb E_{X\sim P_X^{\mathrm{tar}}}
\left[\mathbb E_0\{Y(1)-Y(0)\mid X\}\right].
\label{eq:causal-contrast}
\end{equation}
We assume that the conditional potential-outcome means are defined on the target covariate support and integrable under \(P_X^{\mathrm{tar}}\). This distribution specifies the covariate profiles over which the causal contrast is averaged. For example, taking \(P_X^{\mathrm{tar}}=P_{0,X}\), the marginal covariate distribution under \(P_0\), gives the \emph{average treatment effect} (ATE). Taking \(P_X^{\mathrm{tar}}\) to be the empirical covariate distribution, which assigns probability \(1/n_0\) to each reference observation's covariate vector, gives the sample version of the ATE: the \emph{sample average conditional treatment effect} (SACTE). We will discuss this estimand in more detail in Section~\ref{sec:fixed-design-sacte}. For now, the estimand takes the form in Equation~\eqref{eq:causal-contrast}, where \(P_X^{\mathrm{tar}}\) represents either a population distribution or an empirical distribution induced by the reference sample, and we generally write \(\Phi_0\).

In this paper, we focus on adjustment-based estimation. Identification of the target estimand requires conditional exchangeability and positivity to hold \(P_X^{\mathrm{tar}}\)-almost surely: for \(a\in\{0,1\}\),
\begin{equation}
\text{\textbf{Conditional exchangeability:}}\quad Y(a)\perp A\mid X.
\label{eq:conditional-exchangeability}
\end{equation}
\begin{equation}
\text{\textbf{Positivity:}}\quad P_0(A=a\mid X=x)>0.
\label{eq:positivity}
\end{equation}

The conditions in Equations~\eqref{eq:consistency}, \eqref{eq:conditional-exchangeability}, and~\eqref{eq:positivity} identify the causal conditional means as observed outcome means, and hence identify the estimand from the observed data:
\begin{equation}
\begin{aligned}
\mu_{0,a}(x)
&:=
\mathbb E_0[Y(a)\mid X=x]
=
\mathbb E_0[Y\mid A=a,X=x],\\
\Phi_0
&=
\mathbb E_{X\sim P_X^{\mathrm{tar}}}
\left[\mu_{0,1}(X)-\mu_{0,0}(X)\right].
\end{aligned}
\label{eq:identified-conditional-mean}
\end{equation}

Under conditional exchangeability and positivity, Equation~\eqref{eq:identified-conditional-mean} also accommodates effects among the treated: choosing \(P_X^{\mathrm{tar}}=P_0(X\mid A=1)\) gives the average treatment effect on the treated (ATT), while using the empirical covariate distribution of the treated reference observations gives its sample average conditional analogue.

\subsection{Empirical Monte Carlo benchmarking}
\label{sec:emcs-framework}

The goal of the EMCS is to understand the behavior of a set of adjustment-based causal estimators \(\mathcal E=\{e_1,\dots,e_E\}\) when targeting \(\Phi_0\) under a specific redrawing scheme from \(P_0\). Estimator behavior is quantified using a collection of performance metrics \(\mathcal Q=\{q_1,\dots,q_Q\}\), such as bias, Monte Carlo standard deviation, or mean squared error (MSE). We assume that the chosen metrics are well defined and finite under each benchmark, with the required moment conditions satisfied.

The \emph{target redrawing scheme} specifies the sample size, which variables are held fixed and which are redrawn, and any dependence across units. Together with the original data-generating mechanism, these choices define the dataset law \(\mathbb P_{0,n}\), which generates \(M\) datasets of size \(n\):
\begin{equation}
\mathcal D_{0,n}^{(m)}
\overset{\mathrm{iid}}{\sim}\mathbb P_{0,n},
\qquad m=1,\dots,M.
\label{eq:target-dataset-draws}
\end{equation}
For example, in the independent and identically distributed (iid) case, each dataset is generated by drawing \(n\) independent units from \(P_0\) (here denoting its marginal distribution of the observed variables \((X,A,Y)\)), where \(n\) may differ from the reference sample size \(n_0\):
\[
\begin{gathered}
(X_i,A_i,Y_i)\overset{\mathrm{iid}}{\sim}P_0,
\qquad i=1,\dots,n,\\
\mathcal D_{0,n}=\{(X_i,A_i,Y_i)\}_{i=1}^n
\sim\mathbb P_{0,n}=(P_0)^n.
\end{gathered}
\]
The redrawing scheme can extend to other scenarios. \(\mathbb P_{0,n}\) may instead impose dependence across units $i$, such as group-correlated errors. The target estimand and inference also guide the shape of \(\mathbb P_{0,n}\). For example, if the target is a sample estimand, a natural choice for \(\mathbb P_{0,n}\) would be to generate datasets where covariates remain those of the reference sample, while \(Y_i\) and possibly treatment are redrawn. Together, the dataset law \(\mathbb P_{0,n}\) and the target causal estimand \(\Phi_0\) defined in Equation~\eqref{eq:causal-contrast} form the \emph{target benchmark}, \((\mathbb P_{0,n},\Phi_0)\), which the EMCS will try to approximate.

We denote by \(q(e;\mathbb P_{0,n},\Phi_0)\) the performance of estimator \(e\) under the target benchmark. If the researcher had access to \(\mathbb P_{0,n}\) and \(\Phi_0\), it would suffice to generate \(M\) datasets of size \(n\),
fit estimator \(e\in\mathcal E\) to each dataset, and approximate its performance by the empirical Monte Carlo summary:
\begin{equation}
\begin{gathered}
\widehat q_M
\left(
e;
\mathcal D_{0,n}^{(1)},\dots,\mathcal D_{0,n}^{(M)};
\Phi_0
\right)
\approx q(e;\mathbb P_{0,n},\Phi_0),\\
\text{with }\mathcal D_{0,n}^{(m)}
\overset{\mathrm{iid}}{\sim}\mathbb P_{0,n},
\qquad m=1,\dots,M.
\end{gathered}
\label{eq:target-monte-carlo-performance}
\end{equation}

However, in practice \(P_0\) is unknown, so the target benchmark \((\mathbb P_{0,n},\Phi_0)\) is not available. In an EMCS, the researcher uses \(\mathcal D_0\) to build approximations of both objects: a Monte Carlo dataset law \(\mathbb P_{MC,n}\) and a specified reference value \(\Phi_{MC}\). The pair \((\mathbb P_{MC,n},\Phi_{MC})\) is the \emph{Monte Carlo benchmark}. Synthetic datasets are drawn from this law:
\begin{equation}
\widetilde{\mathcal D}_{MC,n}^{(m)}
\overset{\mathrm{iid}}{\sim}\mathbb P_{MC,n},
\qquad m=1,\dots,M.
\label{eq:synthetic-dataset-draws}
\end{equation}
In this paper, we focus on EMCS that fit conditional outcome models. We use \emph{simulator} to denote the complete stochastic outcome-generating mechanism. Together with the chosen redrawing scheme, the fitted model induces the Monte Carlo dataset law \(\mathbb P_{MC,n}\). Applying Equation~\eqref{eq:identified-conditional-mean} with the conditional means of the same fitted model gives the reference value \(\Phi_{MC}\). We generate synthetic datasets from the Monte Carlo benchmark, fit each estimator on these datasets, and evaluate it against \(\Phi_{MC}\). For each \(e\in\mathcal E\), we would like the resulting empirical performance summary to approximate the summary obtained under the target benchmark:
\begin{equation}
\begin{gathered}
\widehat q_M(e;\widetilde{\mathcal D}_{MC,n}^{(1)},\ldots,\widetilde{\mathcal D}_{MC,n}^{(M)};\Phi_{MC})
\approx\widehat q_M(e;\mathcal D_{0,n}^{(1)},\ldots,\mathcal D_{0,n}^{(M)};\Phi_0),\\
\widetilde{\mathcal D}_{MC,n}^{(m)}\overset{\mathrm{iid}}{\sim}\mathbb P_{MC,n},
\quad
\mathcal D_{0,n}^{(m)}\overset{\mathrm{iid}}{\sim}\mathbb P_{0,n},
\quad m=1,\dots,M.
\end{gathered}
\label{eq:performance-fidelity}
\end{equation}

A particularly interesting use of EMCS is to rank estimators according to some performance metric, for instance based on the MSE. This allows researchers to compare candidate estimators of \(\Phi_0\) when there is model uncertainty. For each metric \(q\in\mathcal Q\), we denote by \(R_q(\mathbb P_{0,n},\Phi_0)\) the ranking of estimators in \(\mathcal E\) according to their performance under the target dataset law \(\mathbb P_{0,n}\). Similarly, \(R_q(\mathbb P_{MC,n},\Phi_{MC})\) denotes the ranking induced by the Monte Carlo strategy. Using MSE as the performance metric, we would ideally have:
\begin{equation}
R_{\mathrm{MSE}}(\mathbb P_{MC,n},\Phi_{MC})
\approx
R_{\mathrm{MSE}}(\mathbb P_{0,n},\Phi_0).
\label{eq:ranking-preservation}
\end{equation}

We refer to the ability of an EMCS to reproduce the target ranking \(R_q(\mathbb P_{0,n},\Phi_0)\) as \emph{ranking fidelity}, the primary criterion studied in the remainder of the paper. Exact agreement between the two rankings corresponds to perfect ranking fidelity. In general, an EMCS is successful when the simulator-induced law \(\mathbb P_{MC,n}\) together with its estimand \(\Phi_{MC}\) reproduces the conclusions about estimator performance that would have been obtained by repeatedly drawing datasets of size \(n\) from the target dataset law \(\mathbb P_{0,n}\). Figure~\ref{fig:conceptual_exercise} summarizes the target and empirical Monte Carlo benchmarking processes.

\begin{figure*}[htbp]
\centering
\begin{tikzpicture}[
    box/.style={
        draw,
        rounded corners,
        align=center,
        font=\fontsize{8pt}{10pt}\selectfont,
        text width=2.4cm,
        minimum height=0.75cm,
        inner sep=2.5pt,
        execute at begin node=\setlength{\parskip}{0pt}
    },
    widebox/.style={box, text width=5.4cm},
    componentbox/.style={box, minimum height=1.4cm},
    finalbox/.style={widebox, minimum height=1.0cm},
    heading/.style={font=\fontsize{8pt}{10pt}\selectfont\bfseries, align=center},
    arrow/.style={-{Latex[length=2.2mm]}, thick}
]

% The two processes have the same structure: inputs, law and target, performance.
\node[heading] at (0,0.8) {Target benchmark};
\node[heading] at (7.8,0.8) {Monte Carlo benchmark};

\node[widebox] (original) at (0,0) {Original causal law \(P_0\)\\ + target covariate distribution \(P_X^{\mathrm{tar}}\)\\ + redrawing scheme};
\node[widebox] (design) at (7.8,0) {Reference dataset \(\mathcal D_0\) + fitted models\\ + target covariate distribution \(P_X^{\mathrm{tar}}\)\\ + redrawing scheme};

\node[componentbox] (redraw) at (-1.7,-1.7) {Target dataset\\ law\\ \(\mathbb P_{0,n}\)};
\node[componentbox] (estimand) at (1.7,-1.7) {Target causal\\ estimand\\ \(\Phi_0\)};
\node[componentbox] (mclaw) at (6.1,-1.7) {Monte Carlo\\ dataset law\\ \(\mathbb P_{MC,n}\)};
\node[componentbox] (mcvalue) at (9.5,-1.7) {Specified\\ reference value\\ \(\Phi_{MC}\)};

\node[finalbox] (targetbenchmark) at (0,-3.4) {Target ranking\\ \(R_{\mathrm{MSE}}(\mathbb P_{0,n},\Phi_0)\)};
\node[finalbox] (mcbenchmark) at (7.8,-3.4) {Monte Carlo ranking\\ \(R_{\mathrm{MSE}}(\mathbb P_{MC,n},\Phi_{MC})\)};

\draw[arrow] ([xshift=-1.7cm]original.south) -- (redraw.north);
\draw[arrow] ([xshift=1.7cm]original.south) -- (estimand.north);
\draw[arrow] ([xshift=-1.7cm]design.south) -- (mclaw.north);
\draw[arrow] ([xshift=1.7cm]design.south) -- (mcvalue.north);

\draw[arrow] (redraw.south) -- ([xshift=-1.7cm]targetbenchmark.north);
\draw[arrow] (estimand.south) -- ([xshift=1.7cm]targetbenchmark.north);
\draw[arrow] (mclaw.south) -- ([xshift=-1.7cm]mcbenchmark.north);
\draw[arrow] (mcvalue.south) -- ([xshift=1.7cm]mcbenchmark.north);

\draw[densely dashed, thick] (targetbenchmark.east)
    -- node[above, font=\fontsize{8pt}{10pt}\selectfont] {\(\approx\)} (mcbenchmark.west);
\end{tikzpicture}

\captionsetup{font=normalsize}
\caption{Conceptual exercise: target and Monte Carlo benchmarks. Ideally, the study should approximate the performance that would be obtained with access to the original objects.}
\label{fig:conceptual_exercise}
\end{figure*}



\subsection{Outcome-model approximation}
\label{sec:outcome-model-approximation}

Most EMCS use \(\mathcal D_0\) to fit an outcome model conditional on covariates and treatment, \(\widehat P_{MC}(Y\mid X,A)\approx P_{0}(Y\mid X,A)\). This model is then employed to construct both \(\mathbb P_{MC,n}\) and \(\Phi_{MC}\). For the same target covariate distribution \(P_X^{\mathrm{tar}}\), replacing the conditional means in Equation~\eqref{eq:identified-conditional-mean} with those of the fitted model gives
\begin{equation}
\begin{aligned}
\widehat\mu_{MC,a}(x)
&:=\mathbb E_{\widehat P_{MC}}[Y\mid X=x,A=a],\\
\Phi_{MC}
&:=\mathbb E_{X\sim P_X^{\mathrm{tar}}}
\left[\widehat\mu_{MC,1}(X)-\widehat\mu_{MC,0}(X)\right].
\end{aligned}
\label{eq:simulator-induced-estimand}
\end{equation}
Therefore, the fitted model affects both the behavior of estimators and the estimand against which they are compared. We can write the fitted conditional mean as
\[
\widehat\mu_{MC,a}(x)=\mu_{0,a}(x)+r_a(x),
\]
where \(r_a\) captures model misspecification and finite-sample estimation error and is fixed once the simulator has been fitted. Subtracting Equation~\eqref{eq:identified-conditional-mean} from Equation~\eqref{eq:simulator-induced-estimand}, with the same target covariate distribution \(P_X^{\mathrm{tar}}\), gives
\begin{equation}
\Phi_{MC}-\Phi_0
=\mathbb E_{X\sim P_X^{\mathrm{tar}}}\left[r_1(X)-r_0(X)\right].
\label{eq:simulator-estimand-shift}
\end{equation}
Errors in the two fitted conditional means can cancel when averaged over the target covariate distribution: \(\Phi_{MC}=\Phi_0\) can hold even when those means are misspecified.

Model alignment bias can arise when an estimator \(e_1\) represents \(\widehat\mu_{MC,a}\) well, whereas another estimator \(e_2\) better approximates \(\mu_{0,a}\). Their MSE ordering can consequently reverse:
\[
\begin{aligned}
\mathrm{MSE}(e_2;\mathbb P_{0,n},\Phi_0)
&<\mathrm{MSE}(e_1;\mathbb P_{0,n},\Phi_0),\\
\mathrm{MSE}(e_2;\mathbb P_{MC,n},\Phi_{MC})
&>\mathrm{MSE}(e_1;\mathbb P_{MC,n},\Phi_{MC}).
\end{aligned}
\]
Such a reversal can occur even when the average contrast in Equation~\eqref{eq:simulator-estimand-shift} is zero, since the fitted outcome structure may alter estimator-specific bias and variance. Preserving the benchmark value therefore does not ensure ranking fidelity as defined in Equation~\eqref{eq:ranking-preservation}.

Matching conditional means alone also does not determine estimator performance. A simulator may fail to reproduce heteroskedasticity; with non-iid errors, even matching each unit's conditional outcome distribution does not ensure that it reproduces the joint outcome law across units. Both discrepancies can affect estimator variance, MSE, and rankings while preserving the reference value.

\subsection{The SACTE estimand and fixed-design benchmarks}
\label{sec:fixed-design-sacte}

Taking the target covariate distribution to be the empirical distribution described in Section~\ref{sec:causal-framework} for a reference sample of size \(n=n_0\) gives the SACTE,
\begin{equation}
\Phi_0^{\mathrm{SACTE}}
=\frac1n\sum_{i=1}^n
\left[\mu_{0,1}(X_i)-\mu_{0,0}(X_i)\right].
\label{eq:sacte}
\end{equation}
This estimand averages conditional treatment effects over the realized covariate profiles, rather than over the population covariate distribution. It therefore requires the conditional outcome means at these profiles but no additional population covariate model. \authorcite{imbensNonparametricEstimationAverage2004} discusses this estimand under the name conditional average treatment effect; we use SACTE to distinguish it from the now-common use of CATE for the pointwise contrast \(\mathbb E[Y(1)-Y(0)\mid X=x]\).

Fixing the reference covariates defines a natural redrawing scheme for this estimand, while treatment assignments can either be redrawn or held fixed, depending on the inferential aim of the simulation. In this paper, we focus on the case where the realized treatment vector is fixed and only outcomes are redrawn conditional on \((X,A)\). The resulting benchmark evaluates performance for the particular covariate--treatment configuration observed in the reference sample. Thus, both the target law \(\mathbb P_{0,n}\) and the simulator law \(\mathbb P_{MC,n}\) will produce the same covariate--treatment configuration, while their outcome mechanisms and corresponding SACTEs may differ. This choice isolates outcome-model approximation bias from approximation of the covariate and treatment distributions and reduces the modeling choices and computational load.

An interesting consequence of targeting the SACTE under this conditional design is that we can retain the observed covariate--treatment pairs without introducing additional bias distortion for propensity-score-based estimation. \authorcite{shawCautionaryNotePlasmode2026} show that the \emph{Sample Treatment} redrawing scheme, which jointly bootstraps $(X,A)$ and generates outcomes, can misrepresent IPW bias relative to the iid target $\mathbb P_{0,n}=(P_0)^n$, even with correctly specified outcome models. This is because its induced law can only generate a single treatment value at unique reference covariate profiles, even when $P_0$ satisfies positivity as in Equation~\eqref{eq:positivity}. Instead, the \emph{Generate Treatment} scheme, which generates treatment from a propensity model, avoids this problem.

This issue does not arise in our fixed-design comparison. For a general IPW estimator, write
\[
\widehat\Phi_{\mathrm{IPW}}=\sum_{i=1}^n w_iY_i,
\]
where \(w_i\) are signed weights determined by covariates, treatment, and estimated propensity scores (see Appendix~\ref{app:ipw} for a description of our specification). In our fixed-design SACTE benchmark, both the target and simulator laws retain the same \((X,A)\). Holding the propensity-fitting folds and randomization fixed gives identical weights under both benchmarks; assume both treatment groups are nonempty and estimated propensity scores lie strictly between zero and one. With conditional outcome means \(\mu_{0,A_i}(X_i)\) and \(\widehat\mu_{MC,A_i}(X_i)\), the difference between conditional biases is
\[
\begin{aligned}
\operatorname{Bias}_{MC}-\operatorname{Bias}_0
&=\sum_{i=1}^n w_i[\widehat\mu_{MC,A_i}(X_i)-\mu_{0,A_i}(X_i)]\\
&\quad-(\Phi_{MC}-\Phi_0)\\
&=\sum_{i=1}^n w_i r_{A_i}(X_i)
-\frac1n\sum_{i=1}^n[r_1(X_i)-r_0(X_i)].
\end{aligned}
\]
Each bias is relative to its corresponding SACTE. The bias discrepancy arises from outcome-model approximation error and tends to zero as the fitted conditional means approach the target means, even with misspecified propensity scores. Because the target law \(\mathbb P_{0,n}\) is itself defined conditionally on the observed \((X,A)\), retaining these values in \(\mathbb P_{MC,n}\) introduces no mismatch in the covariate--treatment configuration. In the Sample Treatment scheme, by contrast, the simulation law replaces the nondegenerate treatment mechanism of an iid target with a degenerate empirical one.
